\clearpage
\VerdierPhysicalPageAnchor{0059}%
\VerdierSemanticAnchor{ch1-1-5-13}
\noindent\textbf{1.5.13.}\quad
本书将使用的 $\mathbb Z^n$-范畴都是严格范畴。因此，我们不打算在一般情形下
讨论 (1.5.12) 所指出的问题。不过，(1.5.12) 的例子表明，在一个范畴
$\mathcal C$ 上确定 $\mathbb Z$-范畴结构格外简单。事实确实如此。
以 $E\text{-}\operatorname{Cat}$ 表示如下 $U$-范畴：其对象是属于 $U$
且装备了一个等价

\VerdierSemanticAnchor{ch1-1-5-13-eq1}
\begin{equation}
  T\colon\mathcal C\longrightarrow\mathcal C
  \tag{1.5.13.1}
\end{equation}
的范畴；其态射是装备了同构

\VerdierSemanticAnchor{ch1-1-5-13-eq2}
\begin{equation}
  m\colon FT\longrightarrow T'F
  \tag{1.5.13.2}
\end{equation}
的函子 $F\colon\mathcal C\to\mathcal C'$。

于是显然有从 $U$-范畴 $\mathbb Z\text{-}\operatorname{Cat}$ 到
$U$-范畴 $E\text{-}\operatorname{Cat}$ 的函子

\VerdierSemanticAnchor{ch1-1-5-13-eq3}
\begin{equation}
  R\colon \mathbb Z\text{-}\operatorname{Cat}
  \longrightarrow E\text{-}\operatorname{Cat}.
  \tag{1.5.13.3}
\end{equation}

\VerdierSemanticAnchor{ch1-1-5-13-action}
函子 $R$ 把 $\mathbb Z$-范畴 $(\mathcal C,\mathbb Z,T,c)$ 映到装备了等价
$T(1)\colon\mathcal C\to\mathcal C$ 的范畴 $\mathcal C$，并把每个
$\mathbb Z$-函子
\[
  (F,m)\colon(\mathcal C,\mathbb Z,T,c)
  \longrightarrow(\mathcal C',\mathbb Z,T',c')
\]
映到装备了同构 $m(1)\colon FT(1)\to T'(1)F$ 的函子 $F$。

\VerdierSemanticAnchor{ch1-1-5-14}
\VerdierSemanticAnchor{ch1-1-5-14-proposition}
\noindent\textbf{命题 1.5.14.}\quad
\emph{函子 $R$ 是范畴的等价。}

\VerdierSemanticAnchor{ch1-1-5-14-proof}
我们只给出证明的若干说明。命题 (1.5.2) 与注 (1.5.8) 表明 $R$ 充分忠实，
因此只需证明 $R$ 本质满。设
$(\mathcal C,T\colon\mathcal C\to\mathcal C)$ 为装备了一个等价的范畴。
取 $T$ 的一个拟逆函子 $T^{-1}\colon\mathcal C\to\mathcal C$，并令

\VerdierSemanticAnchor{ch1-1-5-14-eq1}
\begin{equation}
  \begin{aligned}
    u&\colon\operatorname{id}_{\mathcal C}\longrightarrow TT^{-1},\\
    v&\colon\operatorname{id}_{\mathcal C}\longrightarrow T^{-1}T
  \end{aligned}
  \tag{1.5.14.1}
\end{equation}
为伴随同构。

对每个整数 $n>0$，函子 $T^n$ 与 $(T^{-1})^n$ 互为伴随，因而有伴随同构

\VerdierSemanticAnchor{ch1-1-5-14-eq2}
\begin{equation}
  \begin{aligned}
    u^{(n)}&\colon\operatorname{id}_{\mathcal C}
      \longrightarrow T^n(T^{-1})^n,\\
    v^{(n)}&\colon\operatorname{id}_{\mathcal C}
      \longrightarrow (T^{-1})^nT^n.
  \end{aligned}
  \tag{1.5.14.2}
\end{equation}
